\chapter{Supplementary Material}\label{app:supplementary}

\section{BLE Acceleration Encoding Range}\label{app:ble_acceleration_range}

The BLE payload described in Section~\ref{sec:wearable_hardware_firmware} stores each motion channel as a
signed 16-bit integer. Since one byte contains 8 bits, two bytes contain
16 bits and can therefore store one signed 16-bit value. The representable
integer range is
\begin{equation}
    -2^{15} \leq n \leq 2^{15}-1,
\end{equation}
which corresponds to
\begin{equation}
    -32768 \leq n \leq 32767.
\end{equation}

For the accelerometer channels, the encoded unit is
\(10^{-3}~\si{\meter\per\second\squared}\) per integer count. The physical
acceleration represented by an encoded integer value \(n_a\) is therefore
\begin{equation}
    a = n_a \times 10^{-3}~\si{\meter\per\second\squared}.
\end{equation}
The maximum positive acceleration that can be represented by this format is
\begin{equation}
    a_{\max} = 32767 \times 10^{-3}
    = 32.767~\si{\meter\per\second\squared}.
\end{equation}
The minimum negative acceleration is
\begin{equation}
    a_{\min} = -32768 \times 10^{-3}
    = -32.768~\si{\meter\per\second\squared}.
\end{equation}

Expressed relative to gravitational acceleration, using
\(g \approx 9.81~\si{\meter\per\second\squared}\), this range is approximately
\begin{equation}
    \frac{32.767}{9.81} \approx 3.34~g.
\end{equation}
Thus, the acceleration channels in the 12-byte BLE payload can represent a
physical acceleration range of approximately
\(-32.768\) to \(+32.767~\si{\meter\per\second\squared}\), equivalent to about
\(-3.34~g\) to \(+3.34~g\). Accelerations outside this range cannot be
represented by this payload format without clipping or a different numerical
encoding.

\section{Android Application Identifier and Build Configuration}\label{app:android_build_config}

The Android logger described in Section~\ref{sec:data_logging_inspection_tools}
was implemented as a native Android application. The application identifier and
build settings used for the evaluated version are listed in
Table~\ref{tab:android_build_config}. These values document which Android
application build was used during the receiver-side logging implementation.

\begin{table}[!htbp]
    \centering
    \small
    \caption{Android logger application identifier and build configuration.}
    \label{tab:android_build_config}
    \begin{tabularx}{\linewidth}{@{}p{0.28\linewidth}X@{}}
        \toprule
        Item & Value or meaning \\
        \midrule
        Package name & \path{com.l1172.bleimuandroid}, the Android application identifier used by the operating system and build tools \\
        \texttt{versionName} & \texttt{1.0.1}, the human-readable application version \\
        \texttt{versionCode} & \texttt{2}, the internal Android version number used for update ordering \\
        Compile SDK & 34, the Android SDK version used to compile the application \\
        Target SDK & 34, the Android platform behaviour level targeted by the application \\
        \bottomrule
    \end{tabularx}
\end{table}

\section{Lag Estimation from Acceleration Magnitude}\label{app:lag_estimation}

The Shimmer/XIAO alignment in Section~\ref{sec:signal_alignment_metrics}
estimated the relative timing offset between the two streams using acceleration
magnitude. For each sensor, the three acceleration channels were first combined
into a single magnitude signal:
\begin{equation}
    \lVert \mathbf{a}(t) \rVert =
    \sqrt{a_x(t)^2 + a_y(t)^2 + a_z(t)^2}.
\end{equation}
This converts the three-axis acceleration into one time series for the Shimmer
and one time series for the XIAO.

The two magnitude signals were then compared over a set of candidate temporal
shifts. For each candidate shift \(\tau\), one signal was shifted relative to
the other and the correlation over the overlapping samples was calculated:
\begin{equation}
    r(\tau) =
    \mathrm{corr}\left(a_{\mathrm{Shimmer}}(t),
    a_{\mathrm{XIAO}}(t+\tau)\right).
\end{equation}
The estimated lag was the shift that produced the highest correlation:
\begin{equation}
    \tau^\ast = \arg\max_{\tau} r(\tau).
\end{equation}

In practical terms, this procedure slides one acceleration-magnitude curve
forwards and backwards in time and selects the offset at which the Shimmer and
XIAO curves are most similar. The resulting lag was then used to trim the two
signals to a shared aligned interval before the channel-level comparison
metrics were calculated.

\section{Bus Recording Handling}\label{app:bus_recording_handling}

Bus recordings were attempted as an additional dynamic public-transport
context during the broader pilot field study. They were not used as the main
Objective~3 classifier dataset because the bus sessions were less controlled
than the retained train recordings. Vehicle vibration, intermittent contact,
and power or stream interruptions made the bus data less suitable as the main
evidence for classifier development. Therefore, the main Objective~3 classifier
dataset was limited to clean labelled train-seated segments.

Explicitly labelled bus seated episodes were nevertheless organised as a
supplementary bus-specific dataset. The files were grouped into
\path{no_pelvic_rotation} and \path{with_pelvic_rotation} classes. Only episodes
with explicit labels were used; other rows that might have corresponded to
static sitting were not inferred retrospectively because the label boundary
was uncertain under strong vehicle motion.

A supplementary train-person/test-person analysis was performed using the
script \path{train_chuyan_meixin_test_ruyang.py}. The model was trained on the
Chuyan and Meixin bus source files and tested on the held-out Ruyang bus source
files. This split was used only as an exploratory bus-context check and was not
used for model selection in the main train-seated Objective~3 analysis. The
split is summarised in Table~\ref{tab:bus_supplementary_split}.

\begin{table}[!htbp]
    \centering
    \footnotesize
    \caption{Supplementary bus train/test split.}
    \label{tab:bus_supplementary_split}
    \begin{tabular*}{\linewidth}{@{\extracolsep{\fill}}llrr@{}}
        \toprule
        Split & Class & Sources & Windows \\
        \midrule
        Train & No pelvic rotation & 6 & 253 \\
        Train & With pelvic rotation & 11 & 96 \\
        Test & No pelvic rotation & 3 & 1396 \\
        Test & With pelvic rotation & 4 & 37 \\
        \bottomrule
    \end{tabular*}
\end{table}

The bus analysis used 2~s windows with a 1~s step and 86 features computed from
the six IMU channels, acceleration and gyroscope magnitudes, simple spectral
features, and inter-axis correlations. Five classical classifiers were tested:
logistic regression (LR), support vector machine (SVM), decision tree (DT),
random forest (RF), and k-nearest neighbours (KNN). The positive class was
\path{with_pelvic_rotation}. The held-out Ruyang window-level results are
reported in Table~\ref{tab:bus_supplementary_results}.

\begin{table}[!htbp]
    \centering
    \scriptsize
    \caption{Supplementary bus held-out Ruyang window-level results.}
    \label{tab:bus_supplementary_results}
    \begin{tabular*}{\linewidth}{@{\extracolsep{\fill}}lrrrrrrr@{}}
        \toprule
        Model & Acc. & Bal. acc. & Prec. & Rec. & Spec. & F1 & FP \\
        \midrule
        LR  & 0.378 & 0.681 & 0.040 & 1.000 & 0.361 & 0.077 & 892 \\
        SVM & 0.044 & 0.509 & 0.026 & 1.000 & 0.019 & 0.051 & 1370 \\
        DT  & 0.666 & 0.828 & 0.072 & 1.000 & 0.657 & 0.134 & 479 \\
        RF  & 0.453 & 0.719 & 0.045 & 1.000 & 0.438 & 0.086 & 784 \\
        KNN & 0.435 & 0.710 & 0.044 & 1.000 & 0.420 & 0.084 & 809 \\
        \bottomrule
    \end{tabular*}
\end{table}

All five classifiers detected the 37 held-out pelvic-rotation windows, but the
precision was very low because many no-pelvic-rotation bus windows were
classified as pelvic rotation. The decision tree gave the highest balanced
accuracy among the tested models, but it still produced 479 false-positive
windows. These supplementary results support the decision to keep the bus
recordings outside the main Objective~3 classifier dataset: bus vibration and
context-specific motion can strongly confound pelvic-rotation detection and may
incorrectly reset the inactivity timer in a feedback system.
